\documentclass[11pt]{article}
\usepackage[margin=1in]{geometry}
\usepackage{amsmath,amssymb,amsthm,mathtools}
\usepackage{enumitem}
\usepackage{hyperref}
\newtheorem{theorem}{Theorem}
\newtheorem{definition}{Definition}
\newtheorem{obligation}{Obligation}
\newtheorem{warning}{Warning}
\newcommand{\AZ}{\mathsf A_Z}
\newcommand{\Kcmp}{\mathsf K_{\rm cmp}}
\newcommand{\Kminus}{\mathsf K^-}
\newcommand{\ThetaMinus}{\Theta^-}
\newcommand{\EF}{E_{\rm EF}}
\newcommand{\Esrc}{E_{\rm src}}
\title{Step 70: RH Membrane Target Setup\\From Formed-Layer Membranes to a Zeta-Carrier Construction Program}
\author{Working note}
\date{}
\begin{document}
\maketitle

\section{Purpose}
This note switches the active target from the general formed-layer membrane framework to an RH-facing paper that tries to instantiate the framework on the completed zeta carrier.  We assume, as a scope convention, that the zeta/prime package is already a formed closure with a native membrane.  The remaining problem is not closure-status.  It is the upgrade from native membrane to layer-dissolving membrane for off-critical-zero probes.

The target forward chain is
\[
\boxed{
\AZ\preceq \Kcmp\preceq \ThetaMinus,
\qquad \operatorname{tr}\ThetaMinus\to0,
\quad \text{plus separating readout and all-six records.}
}
\]
Here \(\AZ\) is the zero-side anti-invariant displacement ledger, \(\Kcmp\) is the completed carrier-side anti-invariant currency, and \(\ThetaMinus\) is the accepted anti-invariant budget.

\section{Closure versus RH}
The foundational closure claim and the RH claim are separated as follows.

\begin{description}[leftmargin=2.2cm]
\item[Claim A.] Zeta/primes admits the Six Birds closure vocabulary \(P_1,\dots,P_6\).  This is taken as a formed-closure premise for the present paper.
\item[Claim A'.] Zeta has a native membrane relative to the probes that constitute zeta as a mathematical object.  This is also taken as a closure-level premise.
\item[Claim B.] The native membrane is adequate for the off-critical-zero dissolving probe family, and the anti-invariant budget collapses on a fixed or exhaustive zero ledger.  This is the RH-facing content.
\end{description}

In the notation of the \(\Xi\)-companion, RH is not the existence of native currency.  It is the adequacy upgrade
\[
\Xi_{C}(D_{\rm off}\mid L_{\rm native})\preceq \Omega,
\qquad \Omega\to0\text{ or controlled,}
\]
combined with the root-composite domination and budget collapse records.

\section{Six construction obligations}
The RH paper is organized around six constructive obligations.

\begin{obligation}[O1: Contractive explicit-formula bridge]
Construct feature maps
\[
W_{\rm cmp}=\begin{bmatrix}W_{\rm pr}\\ W_{\infty}\\ W_{\rm pole}\\ W_{\rm tail}\end{bmatrix}
\]
and a zero-side anti-invariant readout \(V_Z\) such that
\[
\boxed{V_Z=T W_{\rm cmp},\qquad \|T\|\le1.}
\]
Equivalently,
\[
\boxed{\AZ=V_Z^*V_Z\preceq W_{\rm cmp}^*W_{\rm cmp}=\Kcmp.}
\]
This is the central analytic-number-theory obligation.  The standard explicit formula is signed; the framework requires a positive feature-map or Douglas-factorization form.
\end{obligation}

\begin{obligation}[O2: Character-source coercive ladder]
Construct upstream-visible source records whose cumulative anti-invariant response-side frame satisfies
\[
F_{\mathcal S,n}\succeq \Lambda_n (\Theta_0^-)^{-1},
\qquad \Lambda_n\to\infty.
\]
Then the anti-invariant carrier budget satisfies
\[
\Kminus_n\preceq \Lambda_n^{-1}\Theta_0^-+\Esrc_n.
\]
Every anti-invariant character block must be covered.  Trace/invariant-only control is not sufficient.
\end{obligation}

\begin{obligation}[O3: Fixed or exhaustive zero ledger]
Provide a completed zero-side ledger \(\AZ\), or finite windows \(\AZ_n\) with a tail bridge
\[
\AZ\preceq \iota_n\AZ_n\iota_n^*+T_n,
\qquad \operatorname{tr}(B_n)+\operatorname{tr}(T_n)\to0.
\]
Without this, shrinking finite-window defects are moving-ledger support evidence only.
\end{obligation}

\begin{obligation}[O4: Exact Hilbert carrier]
Specify the exact completed Hilbert carrier \(H_Z\), audit operator \(C_Z\), native anti-invariant probe family \(L^-
\), dissolving zero probes \(D_{\rm off}\), null-mode quotient, and packaging maps.  Candidate families include adelic, de Branges, Selberg/trace, or a fresh operator construction.
\end{obligation}

\begin{obligation}[O5: Adequacy of native anti-invariant probes]
Prove that the off-critical-zero dissolving probes are visible through the native anti-invariant family, with controlled residual:
\[
\Xi_{C_Z}(D_{\rm off}\mid L^-)\preceq \Omega_Z.
\]
The scalar readout \(\operatorname{Re}(s)-1/2\) separates the critical line, but adequacy is the operator-level condition.
\end{obligation}

\begin{obligation}[O6: All-six channel record]
Populate the records
\[
P_1:\text{ rewrite/gauge},\quad
P_2:\text{ feasibility/null legality},\quad
P_3:\text{ route/duality},\quad
P_4:\text{ staging/refinement},\quad
P_5:\text{ packaging},\quad
P_6:\text{ audit/currency}.
\]
Each channel must carry a defect/status/nonclaim record.
\end{obligation}

\section{Exact confinement theorem specialized to RH}
Combining previous framework theorems gives the following RH-target theorem schema.

\begin{theorem}[RH membrane implication]
Assume O1--O6.  Suppose that for a sequence of strict extensions \(n\),
\[
\AZ\preceq (1+t_n)(\Lambda_n^{-1}\Theta_0^-+\Esrc_n)+(1+t_n^{-1})\EF_n
\]
for a fixed completed zero ledger \(\AZ\), or for an exhaustive ledger with vanishing tail.  If
\[
\operatorname{tr}(\Lambda_n^{-1}\Theta_0^-+\Esrc_n)\to0,
\qquad
\operatorname{tr}\EF_n\to0,
\]
then \(\AZ=0\).  Since the anti-invariant readout
\[
\psi_-(s)=\operatorname{Re}(s)-\frac12
\]
separates the fixed locus of \(J(s)=1-\overline{s}\), every visible zero in the completed ledger lies on
\[
\operatorname{Re}(s)=\frac12.
\]
\end{theorem}

\section{Optimized obstruction budget}
Let
\[
a_n=\operatorname{tr}(\Lambda_n^{-1}\Theta_0^-+\Esrc_n),
\qquad
b_n=\operatorname{tr}\EF_n.
\]
Then
\[
\inf_{t>0}\operatorname{tr}\left[(1+t)(\Lambda_n^{-1}\Theta_0^-+\Esrc_n)+(1+t^{-1})\EF_n\right]
=(\sqrt{a_n}+\sqrt{b_n})^2.
\]
Thus the exact-confinement objective is reduced to forcing both \(a_n\to0\) and \(b_n\to0\) on a fixed/exhaustive zero ledger.

\section{Work plan}
The paper should proceed in three phases.

\subsection*{Phase I: carrier choice and ledger declarations}
Choose a carrier candidate and write the all-six records without claiming positivity.  This is mostly framework-instantiation work.

\subsection*{Phase II: contractive bridge search}
Attempt to express the completed explicit formula as a Douglas factorization.  The immediate deliverable is not a proof of RH but a precise decomposition into positive feature maps and residual defects.

\subsection*{Phase III: coercive source ladder and fixed/exhaustive squeeze}
Construct or rule out natural source ladders and ledger tails.  This is where analytic number theory enters heavily.

\section{Nonclaims}
This setup does not prove RH.  It does not assert that Weil positivity already gives the Douglas factorization.  It does not assert that any known carrier supplies \(\Lambda_n\to\infty\).  It does not treat finite zero windows as exact unless a tail/exhaustivity bridge is supplied.  It assumes zeta/primes is a formed closure, but it does not assume the native membrane is adequate for off-critical-zero dissolving probes.

\end{document}
